% Logement et natalité -- diaporama du bilan au 22 octobre 2026
% Compiler avec LuaLaTeX, deux passes :  lualatex bilan_logement_natalite.tex
% Thème : beamerthemeTDTE.sty (copie de celui du cours d'Angers 2026) et logos/.
\documentclass[10pt,aspectratio=169]{beamer}

\usepackage[french]{babel}
\frenchsetup{AutoSpacePunctuation=false}
\usepackage{csquotes}
\usepackage{booktabs}
\usepackage{amsmath}
\usepackage{pgfplots}
\pgfplotsset{compat=1.18}
\usepackage{beamerthemeTDTE}
\usetikzlibrary{babel}

% couleurs des facteurs de la baisse de l'ICF (celles de la synthèse Word)
\definecolor{fcouple}{HTML}{2A78D6}
\definecolor{flogement}{HTML}{EB6834}
\definecolor{fpf}{HTML}{8DD7BC}
\definecolor{fage}{HTML}{EDA100}
\definecolor{freste}{HTML}{D9D9D6}
% verdicts
\definecolor{tdteok}{HTML}{2C7A4B}
\definecolor{tdtepart}{HTML}{B07206}
\newcommand{\vok}[1]{{\color{tdteok}\bfseries #1}}
\newcommand{\vpart}[1]{{\color{tdtepart}\bfseries #1}}
\newcommand{\vno}[1]{{\color{tdtered}\bfseries #1}}

% style commun des graphiques pgfplots
\pgfplotsset{tdte/.style={
  width=\linewidth, height=5.2cm, axis lines=left, axis line style={tdtegrey!60},
  tick style={tdtegrey!60}, tick label style={font=\tiny,color=tdtegrey},
  label style={font=\scriptsize,color=tdteink}, title style={font=\small\bfseries,color=tdtenavy},
  grid style={tdtelight}, legend style={font=\tiny,draw=none,fill=none}, legend cell align=left,
  scaled ticks=false, /pgf/number format/use comma}}
% barres empilées des causes de la baisse de l'ICF
\pgfplotsset{icf/.style={tdte, xbar stacked, bar width=0.45cm, height=5.3cm, xmin=0, xmax=104,
  xmajorgrids, xtick={0,25,50,75,100}, y=0.78cm, enlarge y limits={abs=0.45cm},
  yticklabel style={font=\scriptsize,color=tdteink,align=right},
  xlabel={part de la baisse nationale de l'ICF (\%)},
  xticklabel={\pgfmathprintnumber[assume math mode=true]{\tick}},
  legend style={at={(0.5,1.03)},anchor=south,legend columns=5,/tikz/every even column/.append style={column sep=6pt}},
  legend image code/.code={\fill[##1] (0cm,-0.08cm) rectangle (0.25cm,0.12cm);}}}

% mêmes barres, en enfants par femme : part de la baisse x 0,44 (ICF de 2,00 à 1,56)
\pgfplotsset{icfpts/.style={icf, xmax=0.47, xtick={0,0.1,0.2,0.3,0.4},
  xlabel={baisse de l'ICF (enfant par femme)},
  extra x ticks={0.44}, extra x tick labels={0{,}44},
  extra x tick style={grid=major, grid style={densely dashed,tdtered},
    tick label style={font=\tiny\bfseries,color=tdtered}}}}

% étiquettes des segments
\tikzset{etiq/.style={font=\tiny\bfseries,anchor=center,inner sep=0pt}}

% légende en carrés pleins
\pgfplotsset{legcarre/.style={legend image code/.code={\fill[##1] (0cm,-0.08cm) rectangle (0.25cm,0.12cm);}}}

\title[Logement et natalité]{Logement et natalité en France}
\subtitle{Bilan du projet SRCV, 2014-2025}
\author[Kévin Genna]{Kévin Genna}
\institute[Chaire TDTE]{Chaire Transitions Démographiques, Transitions Économiques}
\date[22 octobre 2026]{Point d'étape au 22 octobre 2026}

\begin{document}

{\setbeamertemplate{footline}{}
\begin{frame}[noframenumbering]
  \titlepage
\end{frame}}

\begin{frame}{Au programme}
  \begin{columns}[T]
    \begin{column}{0.44\textwidth}
      \small\tableofcontents[hideallsubsections]
    \end{column}
    \begin{column}{0.52\textwidth}
      \begin{callout}{En bref}
        \begin{itemize}\setlength{\itemsep}{2pt}
          \item ICF $-$22\,\% de 2014 à 2025 ; les facteurs mesurés en couvrent les deux tiers.
          \item Vie en couple : un tiers de la baisse.
          \item Petites surfaces : de 4 à 19\,\% ; une contrainte à la marge d'une pièce.
          \item Recensement : résultats confirmés ; moins de premières naissances en ville, faute de pièces.
          \item Politique familiale : de l'ordre de 18\,\%, imprécis.
        \end{itemize}
      \end{callout}
    \end{column}
  \end{columns}
\end{frame}

%%%%%%%%%%%%%%%%%%%%%%%%%%%%%%%%%%%%%%%%%%%%%%%%%%%%%%%%%%%%%%%%%%%%%%%%%%%
\section{La baisse de la fécondité et ses causes}

\begin{frame}{Ce qui explique la baisse de l'ICF entre 2014 et 2025}
  \begin{columns}[T]
    \begin{column}{0.235\textwidth}\kpi[tdtered]{$-$22\,\%}{ICF passé de 2,00 à 1,56 enfant par femme}\end{column}
    \begin{column}{0.235\textwidth}\kpi{34\,\%}{de la baisse viennent du recul de la vie en couple}\end{column}
    \begin{column}{0.235\textwidth}\kpi[tdteorange]{4 à 19\,\%}{pour les petites surfaces, selon la forme retenue}\end{column}
    \begin{column}{0.235\textwidth}\kpi[tdtegrey]{33\,\%}{restent inexpliqués : à situation donnée, moins d'enfants}\end{column}
  \end{columns}
  \vspace{0.4cm}
  \begin{aere}
    \item Aucun facteur n'explique la baisse à lui seul ; ensemble, les facteurs mesurés en couvrent environ \alert{les deux tiers}.
    \item La vie en couple recule à âge donné, surtout à 25-29 ans (de 58,3 à 53,4\,\% des femmes de 15 à 49 ans), dans tous les milieux.
    \item Les parts sont des bornes hautes : elles lisent comme causales des associations, et ne s'additionnent qu'approximativement.
  \end{aere}
\end{frame}

\begin{frame}{La part de chaque facteur}
  \centering\small
  \begin{tabular}{@{}lrrp{0.43\textwidth}@{}}
    \toprule
    Facteur & Part & Fourchette & Lecture \\
    \midrule
    Recul de la vie en couple & 34\,\% & 33 à 34\,\% & Moins de femmes en couple à âge donné \\
    Logement, petites surfaces & 12\,\% & 4 à 19\,\% & IC à 95\,\% de 0 à 33\,\% ; en partie de l'anticipation \\
    Logement, statuts d'occupation & $-$10\,\% & & Recul des propriétaires non accédants, les moins féconds \\
    Politique familiale (barème) & 18\,\% & $-$14 à 52\,\% & Non significatif ; presque entièrement la PAJE \\
    Âge des couples & 4\,\% & 2 à 4\,\% & Les couples ont vieilli d'un demi-an \\
    Prix au m², revenu, ville & $\approx$\,0 & & Aucun effet à surface donnée \\
    Non expliqué & 33\,\% & & Moins de naissances à situation donnée \\
    \bottomrule
  \end{tabular}
  \source{Insee, SRCV 2014-2025 ; OpenFisca-France ; calculs de l'auteur.}
\end{frame}

\begin{frame}{La décomposition par quartile de niveau de vie}
  \begin{columns}[c]
    \begin{column}{0.64\textwidth}
      \begin{tikzpicture}
        \begin{axis}[icf, symbolic y coords={Q4,Q3,Q2,Q1,Ens}, ytick={Q4,Q3,Q2,Q1,Ens},
            yticklabels={{Q4 (aisés)},Q3,Q2,{Q1 (modestes)},Ensemble}]
          \addplot[fill=fcouple,draw=none] coordinates {(33.8000,Ens) (66.0000,Q1) (30.3000,Q2) (33.0000,Q3) (26.1000,Q4)};
          \addplot[fill=flogement,draw=none] coordinates {(12.2000,Ens) (13.6000,Q1) (14.5000,Q2) (3.1000,Q3) (22.8000,Q4)};
          \addplot[fill=fpf,draw=none] coordinates {(17.5000,Ens) (7.3000,Q1) (14.9000,Q2) (32.5000,Q3) (10.8000,Q4)};
          \addplot[fill=fage,draw=none] coordinates {(4.7000,Ens) (0.3000,Q1) (3.4000,Q2) (2.1000,Q3) (9.5000,Q4)};
          \addplot[fill=freste,draw=none] coordinates {(31.8000,Ens) (12.8000,Q1) (36.9000,Q2) (29.3000,Q3) (30.8000,Q4)};
          \legend{Vie en couple, Petites surfaces, Politique familiale, Âge des couples, Non expliqué}
          \coordinate (e1) at (axis cs:16.9000,Ens);
          \coordinate (e2) at (axis cs:39.9000,Ens);
          \coordinate (e3) at (axis cs:54.7500,Ens);
          \coordinate (e4) at (axis cs:84.1000,Ens);
          \coordinate (e5) at (axis cs:33.0000,Q1);
          \coordinate (e6) at (axis cs:72.8000,Q1);
          \coordinate (e7) at (axis cs:83.2500,Q1);
          \coordinate (e8) at (axis cs:93.6000,Q1);
          \coordinate (e9) at (axis cs:15.1500,Q2);
          \coordinate (e10) at (axis cs:37.5500,Q2);
          \coordinate (e11) at (axis cs:52.2500,Q2);
          \coordinate (e12) at (axis cs:81.5500,Q2);
          \coordinate (e13) at (axis cs:16.5000,Q3);
          \coordinate (e14) at (axis cs:52.3500,Q3);
          \coordinate (e15) at (axis cs:85.3500,Q3);
          \coordinate (e16) at (axis cs:13.0500,Q4);
          \coordinate (e17) at (axis cs:37.5000,Q4);
          \coordinate (e18) at (axis cs:54.3000,Q4);
          \coordinate (e19) at (axis cs:64.4500,Q4);
          \coordinate (e20) at (axis cs:84.6000,Q4);
        \end{axis}
        \node[etiq,text=white] at (e1) {34};
        \node[etiq,text=black] at (e2) {12};
        \node[etiq,text=black] at (e3) {18};
        \node[etiq,text=black] at (e4) {32};
        \node[etiq,text=white] at (e5) {66};
        \node[etiq,text=black] at (e6) {14};
        \node[etiq,text=black] at (e7) {7};
        \node[etiq,text=black] at (e8) {13};
        \node[etiq,text=white] at (e9) {30};
        \node[etiq,text=black] at (e10) {14};
        \node[etiq,text=black] at (e11) {15};
        \node[etiq,text=black] at (e12) {37};
        \node[etiq,text=white] at (e13) {33};
        \node[etiq,text=black] at (e14) {32};
        \node[etiq,text=black] at (e15) {29};
        \node[etiq,text=white] at (e16) {26};
        \node[etiq,text=black] at (e17) {23};
        \node[etiq,text=black] at (e18) {11};
        \node[etiq,text=black] at (e19) {10};
        \node[etiq,text=black] at (e20) {31};
      \end{tikzpicture}
    \end{column}
    \begin{column}{0.33\textwidth}
      \small
      \begin{aere}
        \item La politique familiale se concentre sur les \alert{classes moyennes} (33\,\% en Q3, 15\,\% en Q2) : ce sont elles qui ont perdu la PAJE.
        \item Les 66\,\% de la vie en couple en Q1 sont surtout un effet de classement : une femme seule n'a qu'un revenu et tombe en Q1.
      \end{aere}
    \end{column}
  \end{columns}
  \source{Insee, SRCV 2014-2025 ; calculs de l'auteur.}
\end{frame}

\begin{frame}{Par quartile, en enfants par femme}
  \begin{columns}[c]
    \begin{column}{0.64\textwidth}
      \begin{tikzpicture}
        \begin{axis}[icfpts, symbolic y coords={Q4,Q3,Q2,Q1,Ens}, ytick={Q4,Q3,Q2,Q1,Ens},
            yticklabels={{Q4 (aisés)},Q3,Q2,{Q1 (modestes)},Ensemble}]
          \addplot[fill=fcouple,draw=none] coordinates {(0.1487,Ens) (0.2904,Q1) (0.1333,Q2) (0.1452,Q3) (0.1148,Q4)};
          \addplot[fill=flogement,draw=none] coordinates {(0.0537,Ens) (0.0598,Q1) (0.0638,Q2) (0.0136,Q3) (0.1003,Q4)};
          \addplot[fill=fpf,draw=none] coordinates {(0.0770,Ens) (0.0321,Q1) (0.0656,Q2) (0.1430,Q3) (0.0475,Q4)};
          \addplot[fill=fage,draw=none] coordinates {(0.0207,Ens) (0.0013,Q1) (0.0150,Q2) (0.0092,Q3) (0.0418,Q4)};
          \addplot[fill=freste,draw=none] coordinates {(0.1399,Ens) (0.0563,Q1) (0.1624,Q2) (0.1289,Q3) (0.1355,Q4)};
          \legend{Vie en couple, Petites surfaces, Politique familiale, Âge des couples, Non expliqué}
          \coordinate (e1) at (axis cs:0.0744,Ens);
          \coordinate (e2) at (axis cs:0.1756,Ens);
          \coordinate (e3) at (axis cs:0.2409,Ens);
          \coordinate (e4) at (axis cs:0.3700,Ens);
          \coordinate (e5) at (axis cs:0.1452,Q1);
          \coordinate (e6) at (axis cs:0.3203,Q1);
          \coordinate (e7) at (axis cs:0.4118,Q1);
          \coordinate (e8) at (axis cs:0.0667,Q2);
          \coordinate (e9) at (axis cs:0.1652,Q2);
          \coordinate (e10) at (axis cs:0.2299,Q2);
          \coordinate (e11) at (axis cs:0.3588,Q2);
          \coordinate (e12) at (axis cs:0.0726,Q3);
          \coordinate (e13) at (axis cs:0.2303,Q3);
          \coordinate (e14) at (axis cs:0.3755,Q3);
          \coordinate (e15) at (axis cs:0.0574,Q4);
          \coordinate (e16) at (axis cs:0.1650,Q4);
          \coordinate (e17) at (axis cs:0.2389,Q4);
          \coordinate (e18) at (axis cs:0.2836,Q4);
          \coordinate (e19) at (axis cs:0.3722,Q4);
        \end{axis}
        \node[etiq,text=white] at (e1) {0,15};
        \node[etiq,text=black] at (e2) {0,05};
        \node[etiq,text=black] at (e3) {0,08};
        \node[etiq,text=black] at (e4) {0,14};
        \node[etiq,text=white] at (e5) {0,29};
        \node[etiq,text=black] at (e6) {0,06};
        \node[etiq,text=black] at (e7) {0,06};
        \node[etiq,text=white] at (e8) {0,13};
        \node[etiq,text=black] at (e9) {0,06};
        \node[etiq,text=black] at (e10) {0,07};
        \node[etiq,text=black] at (e11) {0,16};
        \node[etiq,text=white] at (e12) {0,15};
        \node[etiq,text=black] at (e13) {0,14};
        \node[etiq,text=black] at (e14) {0,13};
        \node[etiq,text=white] at (e15) {0,11};
        \node[etiq,text=black] at (e16) {0,10};
        \node[etiq,text=black] at (e17) {0,05};
        \node[etiq,text=black] at (e18) {0,04};
        \node[etiq,text=black] at (e19) {0,14};
      \end{tikzpicture}
    \end{column}
    \begin{column}{0.33\textwidth}
      \small
      \begin{itemize}\setlength{\itemsep}{0.4em}
        \item \textbf{France entière}, l'ICF perd 0,44 enfant par femme : 0,15 pour la vie en couple, 0,05 pour les petites surfaces (de 0,02 à 0,08), 0,08 pour la politique familiale, 0,02 pour l'âge, 0,14 non expliqué.
      \end{itemize}
      \vspace{0.15cm}
      \begin{callout}[tdtenavy]{Lire un quartile}
        \footnotesize En Q4, la vie en couple équivaut à 0,11 enfant par femme, au niveau de fécondité national. Ce n'est pas la contribution du Q4 à la baisse de l'ICF de la France.
      \end{callout}
    \end{column}
  \end{columns}
  \source{Insee, SRCV 2014-2025 ; Insee (2026) pour l'ICF ; parts de la diapositive précédente multipliées par 0,44.}
\end{frame}

\begin{frame}{La décomposition par diplôme de la femme}
  \begin{columns}[c]
    \begin{column}{0.64\textwidth}
      \begin{tikzpicture}
        \begin{axis}[icf, symbolic y coords={CAP,Bac,Bac2,Sup,Ens}, ytick={CAP,Bac,Bac2,Sup,Ens},
            yticklabels={{CAP-BEP ou moins},Bac,Bac+2,{Bac+3 ou plus},{Ensemble (25-49 ans)}}]
          \addplot[fill=fcouple,draw=none] coordinates {(30.3000,Ens) (28.6000,Sup) (49.5000,Bac2) (24.9000,Bac) (24.9000,CAP)};
          \addplot[fill=flogement,draw=none] coordinates {(11.7000,Ens) (18.0000,Sup) (0.0000,Bac2) (1.0000,Bac) (15.9000,CAP)};
          \addplot[fill=fpf,draw=none] coordinates {(17.8000,Ens) (16.9000,Sup) (18.6000,Bac2) (19.5000,Bac) (17.6000,CAP)};
          \addplot[fill=fage,draw=none] coordinates {(6.7000,Ens) (8.9000,Sup) (10.6000,Bac2) (1.8000,Bac) (3.8000,CAP)};
          \addplot[fill=freste,draw=none] coordinates {(33.4000,Ens) (27.7000,Sup) (26.2000,Bac2) (52.7000,Bac) (37.8000,CAP)};
          \legend{Vie en couple, Petites surfaces, Politique familiale, Âge des couples, Non expliqué}
          \coordinate (e1) at (axis cs:15.1500,Ens);
          \coordinate (e2) at (axis cs:36.1500,Ens);
          \coordinate (e3) at (axis cs:50.9000,Ens);
          \coordinate (e4) at (axis cs:63.1500,Ens);
          \coordinate (e5) at (axis cs:83.2000,Ens);
          \coordinate (e6) at (axis cs:14.3000,Sup);
          \coordinate (e7) at (axis cs:37.6000,Sup);
          \coordinate (e8) at (axis cs:55.0500,Sup);
          \coordinate (e9) at (axis cs:67.9500,Sup);
          \coordinate (e10) at (axis cs:86.2500,Sup);
          \coordinate (e11) at (axis cs:24.7500,Bac2);
          \coordinate (e12) at (axis cs:58.8000,Bac2);
          \coordinate (e13) at (axis cs:73.4000,Bac2);
          \coordinate (e14) at (axis cs:91.8000,Bac2);
          \coordinate (e15) at (axis cs:12.4500,Bac);
          \coordinate (e16) at (axis cs:35.6500,Bac);
          \coordinate (e17) at (axis cs:73.5500,Bac);
          \coordinate (e18) at (axis cs:12.4500,CAP);
          \coordinate (e19) at (axis cs:32.8500,CAP);
          \coordinate (e20) at (axis cs:49.6000,CAP);
          \coordinate (e21) at (axis cs:81.1000,CAP);
        \end{axis}
        \node[etiq,text=white] at (e1) {30};
        \node[etiq,text=black] at (e2) {12};
        \node[etiq,text=black] at (e3) {18};
        \node[etiq,text=black] at (e4) {7};
        \node[etiq,text=black] at (e5) {33};
        \node[etiq,text=white] at (e6) {29};
        \node[etiq,text=black] at (e7) {18};
        \node[etiq,text=black] at (e8) {17};
        \node[etiq,text=black] at (e9) {9};
        \node[etiq,text=black] at (e10) {28};
        \node[etiq,text=white] at (e11) {50};
        \node[etiq,text=black] at (e12) {19};
        \node[etiq,text=black] at (e13) {11};
        \node[etiq,text=black] at (e14) {26};
        \node[etiq,text=white] at (e15) {25};
        \node[etiq,text=black] at (e16) {20};
        \node[etiq,text=black] at (e17) {53};
        \node[etiq,text=white] at (e18) {25};
        \node[etiq,text=black] at (e19) {16};
        \node[etiq,text=black] at (e20) {18};
        \node[etiq,text=black] at (e21) {38};
      \end{tikzpicture}
    \end{column}
    \begin{column}{0.33\textwidth}
      \small
      \begin{aere}
        \item Le recul de la vie en couple touche \alert{tous les niveaux de diplôme}, sans concentration chez les moins diplômées.
        \item La politique familiale pèse de façon semblable dans chaque groupe (17 à 20\,\%).
        \item Chez les bac+2, la contribution des surfaces est légèrement négative ($-$5\,\%), non représentée.
      \end{aere}
    \end{column}
  \end{columns}
  \source{Insee, SRCV 2014-2025, femmes de 25 à 49 ans ; calculs de l'auteur.}
\end{frame}

\begin{frame}{Par diplôme, en enfants par femme}
  \begin{columns}[c]
    \begin{column}{0.64\textwidth}
      \begin{tikzpicture}
        \begin{axis}[icfpts, symbolic y coords={CAP,Bac,Bac2,Sup,Ens}, ytick={CAP,Bac,Bac2,Sup,Ens},
            yticklabels={{CAP-BEP ou moins},Bac,Bac+2,{Bac+3 ou plus},{Ensemble (25-49 ans)}}]
          \addplot[fill=fcouple,draw=none] coordinates {(0.1333,Ens) (0.1258,Sup) (0.2178,Bac2) (0.1096,Bac) (0.1096,CAP)};
          \addplot[fill=flogement,draw=none] coordinates {(0.0515,Ens) (0.0792,Sup) (0.0000,Bac2) (0.0044,Bac) (0.0700,CAP)};
          \addplot[fill=fpf,draw=none] coordinates {(0.0783,Ens) (0.0744,Sup) (0.0818,Bac2) (0.0858,Bac) (0.0774,CAP)};
          \addplot[fill=fage,draw=none] coordinates {(0.0295,Ens) (0.0392,Sup) (0.0466,Bac2) (0.0079,Bac) (0.0167,CAP)};
          \addplot[fill=freste,draw=none] coordinates {(0.1470,Ens) (0.1219,Sup) (0.1153,Bac2) (0.2319,Bac) (0.1663,CAP)};
          \legend{Vie en couple, Petites surfaces, Politique familiale, Âge des couples, Non expliqué}
          \coordinate (e1) at (axis cs:0.0667,Ens);
          \coordinate (e2) at (axis cs:0.1591,Ens);
          \coordinate (e3) at (axis cs:0.2240,Ens);
          \coordinate (e4) at (axis cs:0.3661,Ens);
          \coordinate (e5) at (axis cs:0.0629,Sup);
          \coordinate (e6) at (axis cs:0.1654,Sup);
          \coordinate (e7) at (axis cs:0.2422,Sup);
          \coordinate (e8) at (axis cs:0.2990,Sup);
          \coordinate (e9) at (axis cs:0.3795,Sup);
          \coordinate (e10) at (axis cs:0.1089,Bac2);
          \coordinate (e11) at (axis cs:0.2587,Bac2);
          \coordinate (e12) at (axis cs:0.3230,Bac2);
          \coordinate (e13) at (axis cs:0.4039,Bac2);
          \coordinate (e14) at (axis cs:0.0548,Bac);
          \coordinate (e15) at (axis cs:0.1569,Bac);
          \coordinate (e16) at (axis cs:0.3236,Bac);
          \coordinate (e17) at (axis cs:0.0548,CAP);
          \coordinate (e18) at (axis cs:0.1445,CAP);
          \coordinate (e19) at (axis cs:0.2182,CAP);
          \coordinate (e20) at (axis cs:0.3568,CAP);
        \end{axis}
        \node[etiq,text=white] at (e1) {0,13};
        \node[etiq,text=black] at (e2) {0,05};
        \node[etiq,text=black] at (e3) {0,08};
        \node[etiq,text=black] at (e4) {0,15};
        \node[etiq,text=white] at (e5) {0,13};
        \node[etiq,text=black] at (e6) {0,08};
        \node[etiq,text=black] at (e7) {0,07};
        \node[etiq,text=black] at (e8) {0,04};
        \node[etiq,text=black] at (e9) {0,12};
        \node[etiq,text=white] at (e10) {0,22};
        \node[etiq,text=black] at (e11) {0,08};
        \node[etiq,text=black] at (e12) {0,05};
        \node[etiq,text=black] at (e13) {0,12};
        \node[etiq,text=white] at (e14) {0,11};
        \node[etiq,text=black] at (e15) {0,09};
        \node[etiq,text=black] at (e16) {0,23};
        \node[etiq,text=white] at (e17) {0,11};
        \node[etiq,text=black] at (e18) {0,07};
        \node[etiq,text=black] at (e19) {0,08};
        \node[etiq,text=black] at (e20) {0,17};
      \end{tikzpicture}
    \end{column}
    \begin{column}{0.33\textwidth}
      \small
      \begin{itemize}\setlength{\itemsep}{0.4em}
        \item Chez les 25-49 ans, la vie en couple équivaut à 0,13 enfant par femme ; la politique familiale à 0,08 dans chaque groupe.
        \item Chez les bac+2, la vie en couple pèse 0,22 ; les surfaces y jouent un peu en sens inverse ($+$0,02, non représenté).
        \item Une barre de groupe se lit comme pour les quartiles.
      \end{itemize}
    \end{column}
  \end{columns}
  \source{Insee, SRCV 2014-2025, femmes de 25 à 49 ans ; parts de la diapositive précédente multipliées par 0,44.}
\end{frame}

%%%%%%%%%%%%%%%%%%%%%%%%%%%%%%%%%%%%%%%%%%%%%%%%%%%%%%%%%%%%%%%%%%%%%%%%%%%
\section{Logement et naissance dans l'enquête SRCV}

\begin{frame}{Les conditions de logement des couples}
  \begin{columns}[T]
    \begin{column}{0.235\textwidth}\kpi{$-$6 m²}{de surface médiane des couples, à moyenne presque inchangée}\end{column}
    \begin{column}{0.235\textwidth}\kpi[tdteorange]{19\,\%}{des couples sans enfant logés sous 50 m², contre 15\,\% en 2014}\end{column}
    \begin{column}{0.235\textwidth}\kpi[tdteorange]{27\,\%}{des couples sans enfant sans place pour un enfant, contre 24\,\%}\end{column}
    \begin{column}{0.235\textwidth}\kpi{$-$6 pts}{de taux d'effort des accédants (de $-$4 à $-$8 selon l'âge)}\end{column}
  \end{columns}
  \vspace{0.3cm}
  \begin{itemize}\setlength{\itemsep}{0.35em}
    \item Les jeunes locataires ont perdu de la place à taux d'effort stable, jusqu'à $-$8\,\% à 35-39 ans.
    \item Le recul se concentre dans le bas de la distribution et dans les grandes agglomérations.
    \item Les parents à qui une pièce manque passent de 5,6 à 10,5\,\%, ce que le recensement ne confirme pas pour 2019-2023.
  \end{itemize}
  \source{Insee, SRCV 2014 et 2025 ; couples dont la femme a de 15 à 49 ans.}
\end{frame}

\begin{frame}{La surface compte, pas son prix ni le revenu}
  \begin{columns}[T]
    \begin{column}{0.30\textwidth}
      \kpi[tdteorange]{+4\,\%}{de probabilité de naissance pour dix m² de plus (+0,34 point sur un taux de 8,2\,\%)}
      \vspace{0.2cm}
      \kpi{1\,519}{naissances observées dans SRCV, 2014-2025, sur 19\,645 couples}
    \end{column}
    \begin{column}{0.66\textwidth}
      \begin{itemize}\setlength{\itemsep}{0.45em}
        \item \textbf{La surface} est associée à la naissance, davantage pour le premier enfant (+0,59 point).
        \item \textbf{Le coût au mètre carré} n'a aucun effet à surface donnée.
        \item \textbf{Le revenu} non plus, s'il est mesuré avant la naissance : le niveau de vie usuel compte le nouveau-né et les prestations qu'il déclenche.
        \item \textbf{Les accédants} gardent un avantage d'environ 5 points au premier enfant.
        \item Chez les couples arrivés depuis moins de trois ans, l'effet est deux fois et demie plus fort : un premier indice d'anticipation.
      \end{itemize}
    \end{column}
  \end{columns}
  \source{Insee, SRCV 2014-2025 ; probits pondérés, effets marginaux moyens.}
\end{frame}

\begin{frame}{Un seuil en pièces pour le premier enfant}
  \begin{columns}[T]
    \begin{column}{0.5\textwidth}
      \begin{tikzpicture}
        \begin{axis}[tdte, ybar, bar width=0.5cm, height=5.4cm, ymin=-15, ymax=6,
            symbolic x coords={a,b,c,d,e,f}, xtick=data,
            xticklabels={{<\,40 m²},{40-49},{50-59},{60-69},{90-119},{120+}},
            x tick label style={font=\scriptsize,color=tdteink}, xticklabel pos=lower,
            ylabel={écart avec 70-89 m² (points)}, ymajorgrids, axis x line*=bottom,
            nodes near coords={\pgfmathprintnumber[fixed,precision=1]{\pgfplotspointmeta}},
            every node near coord/.append style={font=\tiny,color=tdteink},
            title={Probabilité de première naissance}, enlarge x limits=0.1,
            extra y ticks={0}, extra y tick labels={}, extra y tick style={grid style={tdtegrey}}]
          \addplot[fill=tdtenavy,draw=none] coordinates {(a,-12.04) (b,-8.25) (c,-6.69) (d,0.20) (e,3.69) (f,2.08)};
        \end{axis}
      \end{tikzpicture}
    \end{column}
    \begin{column}{0.47\textwidth}
      \small
      \begin{itemize}\setlength{\itemsep}{0.5em}
        \item La relation est une \alert{marche} entre 45 et 60 m², concentrée sur le premier enfant.
        \item Selon la \textbf{norme de peuplement} de l'Insee, ne pas avoir la place d'un enfant abaisse la probabilité de première naissance de \textbf{10 points} ($-$73\,\% du taux).
        \item À place donnée, la surface ne compte plus : il faut un \textbf{trois-pièces}. Un quart des couples sans enfant ne l'ont pas.
      \end{itemize}
    \end{column}
  \end{columns}
  \source{Insee, SRCV 2014-2025 ; tranches de surface en m², couples sans enfant ; probit pondéré.}
\end{frame}

\begin{frame}{Contrainte ou choix, trois tests d'ordre temporel}
  \begin{columns}[T]
    \begin{column}{0.315\textwidth}
      \card{\faHourglassHalf}{Horizon long}{Sur quatre ans, la surface ne prédit pas les naissances suivantes des couples installés.}
    \end{column}
    \begin{column}{0.315\textwidth}
      \card{\faHome}{Place en pièces}{Sans la place d'un enfant, un tiers de naissances en moins ensuite ($-$35\,\%, $p = 0{,}01$).}
    \end{column}
    \begin{column}{0.315\textwidth}
      \card{\faTruck}{Ordre des événements}{Après un agrandissement, deux fois plus de premières naissances ; l'inverse est rare.}
    \end{column}
  \end{columns}
  \vspace{0.3cm}
  \begin{callout}{Verdict}
    Une partie de l'association est de l'anticipation, puisqu'on s'agrandit avant l'enfant. L'autre, environ la moitié de l'écart observé en coupe, est une \textbf{contrainte d'indivisibilité}, la pièce qui manque.
  \end{callout}
\end{frame}

\begin{frame}{Le modèle face aux données}
  \centering\small
  \begin{tabular}{@{}lll@{}}
    \toprule
    Prédiction & Verdict & Détail \\
    \midrule
    P1\enspace La surface élève la natalité & \vok{Vérifiée} & +4\,\% du taux pour dix m² \\
    P2\enspace Le prix ne joue pas à surface donnée & \vok{Vérifiée} & Coût au m² sans effet \\
    P3\enspace Le revenu n'agit que par la surface & \vok{Vérifiée} & Revenu mesuré avant la naissance \\
    P4\enspace La surface prédit les naissances à venir & \vpart{En partie} & Pas la surface, mais la place ($-$35\,\%) \\
    P5\enspace Concavité, modestes plus sensibles & \vpart{En partie} & Seuil en pièces ; pas de gradient de revenu \\
    P6\enspace Moins de naissances en ville & \vpart{En partie} & Premières naissances, faute de pièces \\
    \bottomrule
  \end{tabular}
  \vspace{0.3cm}

  \raggedright\small Calibré, le canal d'espace pur est faible : une élasticité d'environ 0,11, contre 0,4 dans les données. À long terme, la population s'ajuste au parc de logements, pas la surface par ménage.
\end{frame}

%%%%%%%%%%%%%%%%%%%%%%%%%%%%%%%%%%%%%%%%%%%%%%%%%%%%%%%%%%%%%%%%%%%%%%%%%%%
\section{Le recensement, contre-épreuve des résultats}

\begin{frame}{Cent fois plus de naissances que dans SRCV}
  \begin{columns}[T]
    \begin{column}{0.235\textwidth}\kpi{2,2 M}{couples du champ dans le recensement de 2019 (enquêtes 2017-2021)}\end{column}
    \begin{column}{0.235\textwidth}\kpi[tdteorange]{183\,000}{naissances dans l'année en 2019, 143\,000 en 2023}\end{column}
    \begin{column}{0.235\textwidth}\kpi{1,8 M}{couples dans le recensement de 2023 (enquêtes 2021-2025)}\end{column}
    \begin{column}{0.235\textwidth}\kpi[tdtegrey]{1\,519}{naissances dans SRCV sur toute la période}\end{column}
  \end{columns}
  \vspace{0.3cm}
  \begin{columns}[T]
    \begin{column}{0.48\textwidth}
      \begin{callout}[tdteok]{Ce qu'il mesure}
        Pièces, surface en tranches, statut d'occupation, ancienneté d'emménagement, composition exacte du ménage.
      \end{callout}
    \end{column}
    \begin{column}{0.48\textwidth}
      \begin{callout}[tdtered]{Ce qu'il ne mesure pas}
        Ni le revenu ni le coût du logement ; pas de suivi des ménages, donc pas de test d'ordre temporel.
      \end{callout}
    \end{column}
  \end{columns}
  \source{Insee, recensements 2019 et 2023, fichiers détail ; femme de 15 à 49 ans, France métropolitaine.}
\end{frame}

\begin{frame}{Le seuil apparaît sans bruit}
  \begin{columns}[T]
    \begin{column}{0.56\textwidth}
      \begin{tikzpicture}
        \begin{axis}[tdte, ybar, bar width=0.22cm, height=5.3cm, ymin=0, ymax=17,
            symbolic x coords={a,b,c,d,e,f,g}, xtick=data,
            xticklabels={{<\,30},{30-39},{40-59},{60-79},{80-99},{100-119},{120+}},
            x tick label style={font=\scriptsize,color=tdteink},
            ylabel={taux de première naissance (\%)}, ymajorgrids, enlarge x limits=0.08,
            title={Couples sans enfant, par surface (m²)},
            legend style={at={(0.5,-0.2)},anchor=north,legend columns=2,/tikz/every even column/.append style={column sep=8pt}},
            legcarre]
          \addplot[fill=tdtenavy,draw=none] coordinates {(a,6.4) (b,6.3) (c,8.9) (d,14.7) (e,14.4) (f,14.3) (g,14.0)};
          \addplot[fill=tdteorange,draw=none] coordinates {(a,6.2) (b,5.2) (c,7.6) (d,12.3) (e,13.2) (f,14.1) (g,14.2)};
          \legend{RP 2019, RP 2023}
        \end{axis}
      \end{tikzpicture}
    \end{column}
    \begin{column}{0.41\textwidth}
      \small
      \begin{itemize}\setlength{\itemsep}{0.5em}
        \item Le taux de première naissance \alert{bondit à 60 m²}, puis ne varie plus : la marche de SRCV, au même endroit.
        \item Comme dans SRCV, 73\,\% des couples sans enfant ont la place d'un enfant.
        \item Sans cette place, la probabilité de première naissance baisse de \textbf{9,3 points} en 2019, 7,3 en 2023 (SRCV : 10,2).
        \item Place et surface gardent chacune un effet propre.
      \end{itemize}
    \end{column}
  \end{columns}
  \source{Insee, recensements 2019 et 2023 ; naissance dans les douze mois précédant la collecte.}
\end{frame}

\begin{frame}{SRCV et recensement, les résultats comparés}
  \centering\footnotesize
  \begin{tabular}{@{}lccc@{}}
    \toprule
    & SRCV 2014-2025 & RP 2019 & RP 2023 \\
    \midrule
    Couples sans enfant sans place pour un enfant & 26,5\,\% & 26,9\,\% & 24,5\,\% \\
    Pas de place, premier enfant, points (\% du taux) & $-$10,2 ($-$73) & $-$9,3 ($-$76) & $-$7,3 ($-$65) \\
    Pas de place, rang 2 et plus & $-$2,1 ($-$34) & $-$1,8 ($-$28) & $-$1,7 ($-$29) \\
    Premier enfant, logement de 40 à 59 m² & $-$8,3 / $-$6,7 & $-$8,4 & $-$7,2 \\
    Premier enfant, couples installés, petit logement & $-$9,1 & $-$7,4 & $-$6,7 \\
    Arrivés récemment, parents avec / sans naissance & 25,8 / 13,5\,\% & 19,7 / 10,5\,\% & 17,3 / 9,2\,\% \\
    Premier enfant, propriétaires contre locataires & +5,2 & +3,9 & +4,0 \\
    Premier enfant, HLM contre parc privé & & +4,9 & +4,7 \\
    \bottomrule
  \end{tabular}
  \vspace{0.2cm}

  \raggedright\small
  Confirmé, souvent à quelques dixièmes de point près. Au premier enfant, c'est le parc locatif privé qui se distingue.
  \source{Probits (SRCV) et modèles linéaires de probabilité (recensement) ; correspondances dans l'annexe C de l'article.}
\end{frame}

\begin{frame}{Le manque de place pèse aussi chez les installés}
  \begin{columns}[T]
    \begin{column}{0.55\textwidth}
      \begin{tikzpicture}
        \begin{axis}[tdte, ybar, bar width=0.38cm, height=5.3cm, ymin=0, ymax=95,
            symbolic x coords={a,b,c,d}, xtick=data,
            xticklabels={{Tous},{Installés\\$\geq$ 2 ans},{Installés\\$\geq$ 5 ans},{Arrivés\\< 2 ans}},
            x tick label style={font=\scriptsize,color=tdteink,align=center},
            ylabel={premier enfant, baisse (\% du taux)}, ymajorgrids, enlarge x limits=0.15,
            nodes near coords, every node near coord/.append style={font=\tiny,color=tdteink},
            legend style={at={(0.5,1.03)},anchor=south,legend columns=2,/tikz/every even column/.append style={column sep=8pt}},
            legcarre]
          \addplot[fill=tdtenavy,draw=none] coordinates {(a,76) (b,68) (c,62) (d,86)};
          \addplot[fill=tdteorange,draw=none] coordinates {(a,65) (b,58) (c,57) (d,77)};
          \legend{RP 2019, RP 2023}
        \end{axis}
      \end{tikzpicture}
    \end{column}
    \begin{column}{0.42\textwidth}
      \small
      \begin{itemize}\setlength{\itemsep}{0.5em}
        \item Installé depuis deux ans, un couple a conçu son enfant de l'année dans ce logement.
        \item L'effet est le plus fort chez les \alert{arrivés récents} ($-$86\,\% du taux) : la trace de l'anticipation.
        \item Il garde plus des \textbf{quatre cinquièmes} de son ampleur chez les couples installés depuis cinq ans ou plus.
        \item Comme dans SRCV, les parents qui viennent d'avoir un enfant sont deux fois plus souvent arrivés récemment.
      \end{itemize}
    \end{column}
  \end{columns}
  \source{Insee, recensements 2019 et 2023 ; effet rapporté au taux de chaque sous-population.}
\end{frame}

\begin{frame}{Moins de premières naissances en ville}
  \begin{columns}[T]
    \begin{column}{0.56\textwidth}
      \begin{tikzpicture}
        \begin{axis}[tdte, ybar, bar width=0.3cm, height=5.3cm, ymin=-5.6, ymax=3,
            symbolic x coords={a,b,c,d,e}, xtick=data,
            xticklabels={{< 20\,000\\hab.},{20\,000 à\\99\,999},{100\,000\\ou plus},{Petite\\couronne},{Paris}},
            x tick label style={font=\tiny,color=tdteink,align=center},
            ylabel={premier enfant, écart (points)}, ymajorgrids, axis x line=bottom,
            enlarge x limits=0.13,
            extra y ticks={0}, extra y tick labels={}, extra y tick style={grid style={tdtegrey,thin}},
            legend style={at={(0.5,1.03)},anchor=south,legend columns=2,/tikz/every even column/.append style={column sep=8pt}},
            legcarre]
          \addplot[fill=tdtenavy,draw=none] coordinates {(a,-0.64) (b,-1.58) (c,-3.37) (d,-1.26) (e,-4.96)};
          \addplot[fill=tdteorange,draw=none] coordinates {(a,0.08) (b,-0.12) (c,-1.46) (d,2.44) (e,-0.04)};
          \legend{Sans contrôle de la place, Place contrôlée}
        \end{axis}
      \end{tikzpicture}
    \end{column}
    \begin{column}{0.41\textwidth}
      \small
      \begin{itemize}\setlength{\itemsep}{0.45em}
        \item La natalité d'ensemble n'est pas plus faible en ville (8,2\,\% à Paris, 7,0\,\% dans les petites communes).
        \item Mais les \alert{premières naissances} y sont plus rares : $-$5,0 points à Paris.
        \item Sept couples sans enfant parisiens sur dix n'ont pas la place d'un enfant, un sur dix dans les petites communes.
        \item À place donnée, l'écart parisien disparaît ($-$1,4 point en 2023).
      \end{itemize}
    \end{column}
  \end{columns}
  \source{Insee, recensement 2019 ; écart avec les communes non découpées en IRIS, à caractéristiques données.}
\end{frame}

\begin{frame}{Évolution entre 2019 et 2023, et limites}
  \begin{columns}[T]
    \begin{column}{0.48\textwidth}
      \begin{callout}[tdtenavy]{Une baisse uniforme}
        \begin{itemize}\setlength{\itemsep}{2pt}
          \item Taux de naissance des couples : $-$6,9\,\%, presque autant dans chaque situation de logement.
          \item Premier enfant : de 14,4 à 12,9\,\% avec la place d'un enfant, de 6,5 à 5,8\,\% sans.
          \item Les couples sans place passent de 26,9 à 24,5\,\% : la recomposition n'explique \textbf{rien} de la baisse.
        \end{itemize}
      \end{callout}
    \end{column}
    \begin{column}{0.48\textwidth}
      \begin{callout}[tdtered]{Limites}
        \begin{itemize}\setlength{\itemsep}{2pt}
          \item Ni panel, ni revenu, ni coût du logement.
          \item Les deux millésimes partagent l'enquête de 2021.
          \item Taille de commune et non unité urbaine.
          \item Le RP 2014 couvrirait la période de l'article.
        \end{itemize}
      \end{callout}
    \end{column}
  \end{columns}
  \vspace{0.15cm}
  \small Le tout est repris dans l'\textbf{annexe C de l'article} (six tableaux).
\end{frame}

%%%%%%%%%%%%%%%%%%%%%%%%%%%%%%%%%%%%%%%%%%%%%%%%%%%%%%%%%%%%%%%%%%%%%%%%%%%
\section{La politique familiale, une étude à part}

\begin{frame}{Ce que rapporte un enfant de plus}
  \begin{columns}[T]
    \begin{column}{0.28\textwidth}
      \kpi{4\,866\,€}{par an pour un enfant de plus, barème de 2014 indexé}
      \vspace{0.2cm}
      \kpi[tdtered]{3\,772\,€}{barème de 2025 : $-$22\,\%, et $-$33\,\% au premier enfant}
    \end{column}
    \begin{column}{0.68\textwidth}
      \footnotesize
      \begin{tabular}{@{}lrrr@{}}
        \toprule
        Instrument & 2014 & 2025 & Part de la baisse \\
        \midrule
        Prestations monétaires & 3\,666\,€ & 2\,635\,€ & 16\,\% \\
        \quad dont PAJE & 2\,211\,€ & 1\,378\,€ & 14\,\% \\
        \quad dont AF et complément familial & 1\,455\,€ & 1\,257\,€ & 2\,\% \\
        Quotient familial & 1\,200\,€ & 1\,137\,€ & 1\,\% \\
        Crèches, assistantes maternelles & \multicolumn{2}{c}{offre en hausse} & non chiffré \\
        \bottomrule
      \end{tabular}
      \vspace{0.3cm}

      \small Le gain d'un enfant de plus est simulé avec \textbf{OpenFisca}, en local, selon le revenu du couple, l'activité des parents et l'âge des enfants déjà nés.
    \end{column}
  \end{columns}
  \source{Insee, SRCV 2025 ; OpenFisca-France, barèmes 2014 et 2025 ; calculs de l'auteur.}
\end{frame}

\begin{frame}{Qui a perdu, et avec quel effet}
  \begin{columns}[T]
    \begin{column}{0.56\textwidth}
      \footnotesize\setlength{\tabcolsep}{4pt}
      \begin{tabular}{@{}lrl@{}}
        \toprule
        Quartile & Recul du gain & Instrument en cause \\
        \midrule
        Q1 (modestes) & $-$512\,€ ($-$11\,\%) & PAJE, en partie \\
        Q2 & $-$1\,583\,€ ($-$25\,\%) & PAJE, quotient familial \\
        Q3 & $-$1\,798\,€ ($-$35\,\%) & PAJE, plafonds abaissés \\
        Q4 (aisés) & $-$483\,€ ($-$14\,\%) & Modulation des AF \\
        \bottomrule
      \end{tabular}
    \end{column}
    \begin{column}{0.40\textwidth}
      \small
      \begin{aere}
        \item Mille euros de gain en plus : +0,29 point de probabilité de naissance, \alert{non significatif} ($p = 0{,}26$).
        \item L'imprécision vient du petit nombre de naissances ; il en faudrait environ 9\,300.
        \item D'où une étude à part, fondée sur l'\textbf{ERFS 2010-2023} : demande à adresser à l'ADISP.
      \end{aere}
    \end{column}
  \end{columns}
  \source{Insee, SRCV 2014-2025 ; OpenFisca-France ; quartiles de niveau de vie des couples.}
\end{frame}

%%%%%%%%%%%%%%%%%%%%%%%%%%%%%%%%%%%%%%%%%%%%%%%%%%%%%%%%%%%%%%%%%%%%%%%%%%%
\section{Données, pistes et suite}

\begin{frame}{Les pistes explorées en octobre}
  \begin{columns}[T]
    \begin{column}{0.235\textwidth}
      \card{\faUsers}{Vie en couple}{À natalité inchangée, son recul explique un tiers de la baisse de l'ICF, dans tous les milieux.}
    \end{column}
    \begin{column}{0.235\textwidth}
      \card{\faCity}{Géographie}{L'effet de composition est nul ; les premières naissances sont plus rares en ville, faute de pièces.}
    \end{column}
    \begin{column}{0.235\textwidth}
      \card{\faCalendar}{Report}{L'âge à la maternité augmente à rythme constant : le report n'explique pas la baisse.}
    \end{column}
    \begin{column}{0.235\textwidth}
      \card{\faKey}{Départ}{36 à 38\,\% des 18-34 ans vivent chez leurs parents, sans évolution ; départ plus tardif à Paris.}
    \end{column}
  \end{columns}
\end{frame}

\begin{frame}{Les anomalies trouvées dans les données}
  \begin{aere}
    \item \textbf{Pondération 2022-2023} faussée pour les ménages avec nouveau-né. Recalée sur les naissances de l'état civil : aucun résultat ne change.
    \item \textbf{Double comptage} : dans les vagues 2022-2025, la moitié des naissances est comptée deux fois ; le niveau de natalité mesuré chute avec la refonte de 2022.
    \item \textbf{Valeurs aberrantes} : surfaces du fichier de panel, 38 mensualités d'emprunt, type de logement qui change de nomenclature sous le même nom, tous corrigés.
    \item \textbf{Code géographique} absent de la nomenclature en 2014 : il s'agit du Sud-Ouest.
  \end{aere}
  \vspace{0.2cm}
  \begin{callout}{À signaler}
    Un courrier à l'ADISP sur l'anomalie de pondération est prêt, en brouillon.
  \end{callout}
\end{frame}

\begin{frame}{Ce qui reste à faire}
  \begin{columns}[T]
    \begin{column}{0.31\textwidth}
      \begin{callout}{Avant diffusion}
        \begin{itemize}\setlength{\itemsep}{2pt}
          \item Fusionner et pousser la branche de travail.
          \item Envoyer le courrier à l'ADISP.
          \item Deux études : l'article sur le logement, la politique familiale à part.
        \end{itemize}
      \end{callout}
    \end{column}
    \begin{column}{0.31\textwidth}
      \begin{callout}[tdtenavy]{Questions de fond}
        \begin{itemize}\setlength{\itemsep}{2pt}
          \item Causalité : seule une variation extérieure de l'offre trancherait.
          \item Vie en couple et départ de chez les parents.
          \item Crédit, rang et pièces hors du modèle.
        \end{itemize}
      \end{callout}
    \end{column}
    \begin{column}{0.31\textwidth}
      \begin{callout}[tdteok]{Données}
        \begin{itemize}\setlength{\itemsep}{2pt}
          \item ERFS 2010-2023 pour la politique familiale (naissances repérables, pas de surface).
          \item RP 2014 pour couvrir la période de l'article.
        \end{itemize}
      \end{callout}
    \end{column}
  \end{columns}
\end{frame}

\begin{frame}{Les documents produits}
  \centering\small
  \begin{tabular}{@{}p{0.5\textwidth}rp{0.33\textwidth}@{}}
    \toprule
    Document & Pages & État \\
    \midrule
    Article, « La surface fait-elle les naissances ? » & 32 & À jour au 22 octobre, avec l'annexe sur le recensement \\
    Étude SRCV, 2006-2025 & 42 & À jour au 21 octobre ; mise de côté \\
    Synthèse Word, 2014-2025 & 9 & 21 octobre \\
    Note de synthèse & 4 & 21 octobre \\
    Note de modélisation & 10 & 1er octobre \\
    Ce diaporama & \inserttotalframenumber & 22 octobre \\
    \bottomrule
  \end{tabular}
\end{frame}

\begin{frame}{À retenir}
  \begin{enumerate}\setlength{\itemsep}{0.7em}
    \item La baisse de la fécondité depuis 2014 tient d'abord au \textbf{recul de la vie en couple} ; un tiers reste inexpliqué.
    \item Le logement en porte une part limitée, de \textbf{4 à 19\,\%} au plus, et la plupart de l'association est un logement choisi en vue d'un enfant.
    \item Une contrainte existe pourtant, \textbf{à la marge d'une pièce} : sans la place d'un enfant, un tiers de naissances en moins dans les années suivantes.
    \item Le recensement le confirme, et la situe surtout dans les grandes villes, où la plupart des couples sans enfant n'ont pas cette pièce.
    \item Une politique du logement qui viserait la natalité porterait sur l'\textbf{offre de logements familiaux}, en nombre de pièces.
  \end{enumerate}
\end{frame}

\end{document}
